\documentclass[11pt,a4paper]{ctexart}
\usepackage[margin=2.5cm]{geometry}
\usepackage{amsmath,amssymb,mathtools}
\usepackage[colorlinks=true,linkcolor=blue,urlcolor=blue]{hyperref}
\usepackage{enumitem}
\setlist{nosep}
\setlength{\parindent}{2em}
\setlength{\parskip}{0.35em}

\title{数学分析阅读记录}
\author{}
\date{}

\begin{document}
\maketitle

\section{含参变量的积分}

% STUDY_RECORD_BEGIN id=analysis-parameter-improper-uniform-boundary
\subsection{参数端点发散与内区间上一致收敛的矛盾}
\label{note:analysis-parameter-improper-uniform-boundary}

\textbf{问题与出处：}含参变量的反常积分。设
\[
 f:[0,+\infty)\times[0,1]\longrightarrow\mathbb R
\]
连续。对 $y=0$，积分 $\int_0^{+\infty}f(x,0)\,dx$ 发散；对每个
$y\in(0,1]$，积分 $\int_0^{+\infty}f(x,y)\,dx$ 收敛。问它是否必定不在
$y\in(0,1)$ 上一致收敛，并希望用函数列在端点处极限不存在的对应结论来解释。

\textbf{简短解释：}对每个截断上限 $R>0$，令
\[
 F_R(y)=\int_0^R f(x,y)\,dx,\qquad y\in[0,1].
\]
每个 $F_R$ 都在 $[0,1]$ 上连续。若原反常积分在 $(0,1)$ 上一致收敛，
则一致 Cauchy 准则与有限区间含参积分关于 $y$ 的连续性允许令
$y\to0^+$，从而推出 $F_R(0)$ 在 $R\to+\infty$ 时满足 Cauchy 准则。
然而
\[
 F_R(0)=\int_0^R f(x,0)\,dx,
\]
这将迫使 $y=0$ 处的反常积分收敛，与题设矛盾。因此原积分不可能在
$(0,1)$ 上一致收敛。这里的对应是把函数列 $f_n(y)$ 换成截断积分族
$F_R(y)$，把 $n\to\infty$ 换成 $R\to+\infty$。

% STUDY_RECORD_END id=analysis-parameter-improper-uniform-boundary

\end{document}
